\section{Discussion and Limitations}
\label{sec:discussion}

\subsection{Interpreting a second scheduler}

The crucial distinction is between forming a batch by waiting and forming one
because the outer worker was already busy. The former can be tuned with a
micro-window; the latter is a consequence of serial service occupancy. If
observed batch size rises while intentional companion probability remains near
zero, backlog coalescing is the parsimonious explanation. Whether that is
beneficial depends on the full latency decomposition. Lower backend time does
not imply lower end-to-end latency when the same policy increases queue time or
withholds prompts from continuous batching.

This distinction also changes how an adaptive policy should be judged. A
low-load bypass can correctly decide never to wait and still emit multi-request
batches after a backend call. Calling such behavior ``adaptive waiting'' would
be misleading; it is adaptive \emph{closure} of backlog. A future gateway that
wants to coordinate rather than compete with \vllm{} may need multiple
concurrent backend calls, token-level feedback, or an engine-native scheduling
interface. Those are different designs from a black-box HTTP batcher.

\subsection{Practical decision rule}

For a gateway operator, the first question should be whether direct \vllm{}
already achieves the offered load within the latency objective. If it does,
outer batching must demonstrate an end-to-end improvement, not merely a larger
batch size. At low $\lambda w$, a fixed window has little opportunity to find
companions and should default to pass-through. At high $\lambda S$, a serial
gateway will coalesce backlog automatically, but operators should inspect queue
growth and tail latency before treating this as a throughput optimization.
The safest default is therefore engine-owned batching with gateway
instrumentation; enable outer coalescing only when paired measurements show
that reduced backend cost exceeds added queueing for the target workload.

\subsection{Threats to validity}

\textbf{Hardware and runtime.} The study uses one consumer GPU under WSL2 and
the Windows display-driver model. Absolute latency, host scheduling, power,
and thermal behavior may differ on native Linux or datacenter GPUs. We treat
arrival drift as a validity variable precisely because the host is part of the
measured system.

\textbf{Model and workload.} \model{} is small enough to run reproducibly on
12\,GB VRAM. Larger models, long contexts, quantization, tensor parallelism,
and speculative decoding can change service time and the balance between
prefill and decode. The primary workload uses short, synthetic prompts,
constant arrivals, and a 64-token output cap. Real traces are bursty and
heterogeneous \citep{wang2024burstgpt}; our result should not be generalized to
them without the planned sensitivity study.

\textbf{Load range.} The evaluated rates are calibration-derived points up to
1.8 requests/s. They test a regime in which direct service accepts all offered
traffic, not maximum sustainable saturation. This is deliberate for testing
the hypothesis, but it does not establish the gateway's overload capacity.

\textbf{Statistics.} Three repetitions support paired mechanism evidence but
only coarse uncertainty estimates. Student-$t$ intervals with two degrees of
freedom are sensitive to any outlying run. We therefore show individual run
points and avoid significance language unsupported by the interval.

\textbf{Implementation boundary.} \adip{} uses one coordinator that awaits a
backend batch. The measured backlog mechanism is specific to this transparent,
serial outer architecture. A pipelined or concurrent gateway could expose
arrivals to \vllm{} sooner, while an engine-integrated policy could access
token-level state. The paper characterizes a common and intentionally simple
design, not every possible gateway.

\textbf{Metric boundary.} End-to-end request latency is the principal outcome.
The completions API does not expose token streaming in this experiment, so we
do not separately report time to first token and time per output token. A
system optimized for interactive streaming may require those objectives.

\subsection{Reproducibility as a result}

The artifact separates exploratory pilot evidence from final claims, freezes
policy parameters before the primary matrix, records every exclusion, and
computes uncertainty across runs. This discipline matters for small-system
research: a negative result with a complete causal trail is more informative
than a large percentage derived from one favorable run. The code, configs,
tests, raw-schema definitions, and analysis scripts are intended to make every
number in the final paper traceable to named run IDs.
